\documentclass[11pt]{article}

\usepackage[T1]{fontenc}
\usepackage{lmodern}
\usepackage{microtype}
\usepackage[a4paper,margin=28mm]{geometry}
\usepackage{amsmath,amssymb,amsthm,mathtools}
\usepackage[numbers,sort&compress]{natbib}
\usepackage[hidelinks]{hyperref}

\hypersetup{
  pdftitle={The third boundary coefficient of even off-diagonally symmetric alternating-sign matrices},
  pdfauthor={Xiaojian Zeng}
}

\allowdisplaybreaks

\newtheorem{theorem}{Theorem}
\newtheorem{lemma}[theorem]{Lemma}
\newtheorem{remark}[theorem]{Remark}

\newcommand{\OSASM}{\operatorname{OSASM}}
\newcommand{\Pf}{\operatorname{Pf}}
\newcommand{\ogf}{\operatorname{ogf}}
\newcommand{\Z}{\mathbb{Z}}
\newcommand{\Q}{\mathbb{Q}}

\title{The third boundary coefficient of even off-diagonally symmetric
alternating-sign matrices}
\author{Xiaojian Zeng\\
\small PBC School of Finance, Tsinghua University, 100190\\
\small \texttt{xiaojianzeng@foxmail.com}}
\date{}

\begin{document}
\maketitle

\begin{abstract}
Let \(X^{\mathrm O}_{2n}(r,t)\) be the generating polynomial of
off-diagonally symmetric alternating-sign matrices of order \(2n\), where
\(r\) records the number of nonzero entries strictly above the diagonal and
\(t\) records the position of the \(1\) in the first row.  Behrend, Fischer
and Koutschan conjectured a closed expression for its third boundary
coefficient.  We prove, for every \(n\geq 2\) and as an identity in
\(\Z[r]\),
\[
 [t^3]X^{\mathrm O}_{2n}(r,t)
 =r\bigl(1+(n-2)r^2\bigr)X^{\mathrm O}_{2n-2}(r,1).
\]
The source Pfaffian reduces the problem to a relation between two adjacent
minors of a skew coefficient matrix.  A lower-triangular binomial
congruence centres its kernel and makes the resulting matrix bipartite with
respect to index parity.  Pfaffian minor summation then annihilates the
exceptional shifted minor and proves the required relation for all sizes.
The result is an exact finite-size boundary identity for the associated
six-vertex model; no thermodynamic limit is taken.
\end{abstract}

\noindent\textbf{Keywords:} alternating-sign matrix; six-vertex model;
Pfaffian; boundary correlation; exact enumeration

\section{Introduction}

Alternating-sign matrices (ASMs) and six-vertex partition functions provide
a well-developed bridge between exact enumeration and integrable statistical
mechanics.  Symmetry classes of ASMs lead to modified boundary geometries and
to determinant or Pfaffian partition functions
\cite{Kuperberg2002,RazumovStroganov2004}.  In particular, Behrend, Fischer
and Koutschan (BFK) introduced a triangular-domain six-vertex construction
for diagonally symmetric ASMs and obtained Pfaffian generating functions for
several statistics \cite{BFK2026}.

The even off-diagonally symmetric specialization has zero diagonal.  BFK
proved a reflection relation and the coefficient formula at the second
admissible boundary position, but stated the corresponding third-position
formula for a general weight \(r\) as a conjecture
\cite[Eq.~(8.33)]{BFK2026}.  Refined enumeration at \(r=1\) follows from
earlier work \cite{RazumovStroganov2004}, while the domino-tiling result of
Lee supplies the \(r=\sqrt 2\) specialization
\cite{Lee2024,BFK2026}.  Kumari's subsequent results resolve the three
OSASM conjectures from the earlier version of BFK, rather than this new
formal-\(r\) identity \cite{Kumari2026,BFK2026}.

Here we prove the formal-\(r\) conjecture in full.  The mechanism has two
parts.  First, a low-\(t\) Laurent expansion of the last column in the BFK
Pfaffian reduces the desired coefficient to two adjacent Pfaffian minors.
Second, a binomial change of basis centres the bulk kernel.  In the centred
basis, entries with equal index parity vanish.  A parity imbalance therefore
kills the only extra term in the Pfaffian minor-summation formula
\cite{IshikawaWakayama1995}.

\section{Model and result}

An ASM of order \(m\) is a matrix with entries in
\(\{-1,0,1\}\) whose nonzero entries alternate in each row and column and
whose row and column sums are all \(1\).  For \(n\geq 1\), let
\[
 \OSASM(2n)=
 \left\{A:\ A\text{ is an ASM of order }2n,\quad
 A=A^{\mathsf T},\quad A_{ii}=0\right\}.
\]
For \(A\in\OSASM(2n)\), define
\[
 R(A)=\sum_{1\leq i<j\leq 2n}|A_{ij}|,\qquad
 T(A)=\text{the column of the unique \(1\) in the first row}.
\]
Thus \(T(A)\in\{2,\ldots,2n\}\).  The refined polynomial is
\begin{equation}
 X^{\mathrm O}_{2n}(r,t)
 =\sum_{A\in\OSASM(2n)}r^{R(A)}t^{T(A)}
 \in\Z[r,t].
 \label{eq:generator}
\end{equation}
Under the BFK bijection, this is a finite triangular-domain six-vertex
partition function with a marked boundary position \cite{BFK2026}.  Both
weights in \eqref{eq:generator} are dimensionless.

\begin{theorem}
\label{thm:main}
For every integer \(n\geq2\),
\begin{equation}
 [t^3]X^{\mathrm O}_{2n}(r,t)
 =r\bigl(1+(n-2)r^2\bigr)X^{\mathrm O}_{2n-2}(r,1)
 \qquad\text{in }\Z[r].
 \label{eq:main}
\end{equation}
\end{theorem}

Equivalently, if \(N_n(\rho,\tau)\) counts the matrices with
\((R,T)=(\rho,\tau)\), and \(C_n(\rho)\) counts them with \(R=\rho\), then
\begin{equation}
 N_n(\rho,3)
 =C_{n-1}(\rho-1)+(n-2)C_{n-1}(\rho-3).
 \label{eq:coefficientwise}
\end{equation}
Equation \eqref{eq:coefficientwise} is an algebraic consequence of the
proof; a direct bijection realizing its two terms is not asserted.

\section{Reduction to an adjacent Pfaffian minor}

Let \(m=2n\) and \(k=n-1\).  Write \(A=(a_{ij})_{i,j\geq0}\) for the skew
coefficient matrix with ordinary generating function
\begin{equation}
 A(u,v)=\sum_{i,j\geq0}a_{ij}u^iv^j
 =\frac{(v-u)(1+u)(1+v)}
 {(1-uv)\bigl((1-ru)(1-rv)-uv\bigr)}.
 \label{eq:A-kernel}
\end{equation}
This is the coefficient kernel in the BFK Pfaffians
\cite[Eqs.~(8.10)--(8.11)]{BFK2026}.  For an even-cardinality index set
\(I\), \(A_I\) denotes the corresponding principal submatrix.  Define
\begin{equation}
 C_0=\Pf A_{\{0,1,\ldots,2k-1\}},\qquad
 C_1=\Pf A_{\{0,1,\ldots,2k-2,2k\}}.
 \label{eq:C0C1}
\end{equation}

\begin{lemma}
\label{lem:reduction}
The BFK Pfaffian formula implies
\begin{align}
 [t^3]X^{\mathrm O}_{2n}(r,t)
 &=r^n\bigl((1-r-r^2)C_0+rC_1\bigr),
 \label{eq:lowt}\\
 X^{\mathrm O}_{2n-2}(r,1)&=r^{n-1}C_0.
 \label{eq:unrefined}
\end{align}
\end{lemma}

\begin{proof}
For even \(m\), the non-Pfaffian term in BFK Eq.~(8.11) is absent.
The interior block of its skew matrix is \eqref{eq:A-kernel}; if
\(b_i(t)\) is the entry in row \(i\) of the last column, then
\begin{equation}
 b_i(t)=[u^i]\,
 \frac{(t-1-rtu)(1+u)}
 {(rt-(t-1)u)\bigl(r-(t-1+r^2)u\bigr)}.
 \label{eq:last-column}
\end{equation}
For \(i\geq1\), expansion of the two denominator factors as geometric
series gives the two lowest Laurent coefficients
\begin{equation}
 [t^{-i-1}]b_i=(-1)^{i+1}r^{-i-2},\qquad
 [t^{-i}]b_i=(-1)^i(i+r+r^2)r^{-i-2}.
 \label{eq:Laurent}
\end{equation}
For the first equality, only the constant term of the numerator and the
summand with geometric degrees \((i,0)\) contribute.  At the next power,
the contributions have degrees \((i,0)\), \((i-1,1)\), and
\((i-1,0)\) from the term linear in \(u\); their sum yields the second
coefficient in \eqref{eq:Laurent}.

Expanding the Pfaffian along its last column, with sign \((-1)^i\) for
row \(i\), shows that only \(i=m-2,m-3\) can affect \([t^3]\).
Equations \eqref{eq:C0C1} and \eqref{eq:Laurent} give
\begin{align}
 [t^{-(m-1)}]\Pf M&=-r^{-m}C_0,\nonumber\\
 [t^{-(m-2)}]\Pf M
 &=r^{-m}(m-2+r+r^2)C_0-r^{-(m-1)}C_1.
 \label{eq:pf-laurent}
\end{align}
Now BFK Eq.~(8.11) reads
\[
 X^{\mathrm O}_{m}(r,t)
 =r^{3m/2}t^{m+1}(t-1)^{-(m-1)}\Pf M.
\]
Since \(m-1\) is odd,
\[
 (t-1)^{-(m-1)}
 =-\bigl(1+(m-1)t+O(t^2)\bigr).
\]
Substitution of \eqref{eq:pf-laurent}, followed by \(m=2n\), proves
\eqref{eq:lowt}.  The specialization of BFK Eq.~(8.10) at order \(m-2\)
is exactly \eqref{eq:unrefined}.
\end{proof}

By Lemma~\ref{lem:reduction}, Theorem~\ref{thm:main} will follow once we
establish
\begin{equation}
 C_1=(1+kr)C_0.
 \label{eq:minor-target}
\end{equation}

\section{Binomial centring of the skew kernel}

Remove the two elementary boundary factors from \eqref{eq:A-kernel}:
\begin{equation}
 B_r(u,v)=
 \frac{v-u}{(1-uv)\bigl((1-ru)(1-rv)-uv\bigr)},\qquad
 A(u,v)=(1+u)(1+v)B_r(u,v).
 \label{eq:B-kernel}
\end{equation}
For a scalar \(c\), introduce the infinite unit lower-triangular binomial
matrix
\begin{equation}
 T(c)_{ij}=
 \begin{cases}
 \binom{i}{j}c^{i-j},&i\geq j,\\
 0,&i<j.
 \end{cases}
 \label{eq:T}
\end{equation}
It satisfies \(T(c)T(d)=T(c+d)\).  If \(K(u,v)\) is the ordinary
generating function of a coefficient matrix \(K\), direct summation in
\eqref{eq:T} gives
\begin{equation}
 \ogf\!\left(T(c)KT(c)^{\mathsf T}\right)
 =\frac{1}{(1-cu)(1-cv)}
 K\!\left(\frac{u}{1-cu},\frac{v}{1-cv}\right).
 \label{eq:congruence-ogf}
\end{equation}

Apply \eqref{eq:congruence-ogf} to \(B_r\) with \(c=-r/2\), and denote
the transformed coefficient matrix by \(E\).  Rational simplification
yields
\begin{equation}
 E(u,v)=
 \frac{v-u}{
 \bigl((1+ru/2)(1+rv/2)-uv\bigr)
 \bigl((1-ru/2)(1-rv/2)-uv\bigr)}.
 \label{eq:E-kernel}
\end{equation}
The two denominator factors in \eqref{eq:E-kernel} are interchanged by
\((u,v)\mapsto(-u,-v)\), whereas the numerator changes sign.  Hence
\begin{equation}
 E(-u,-v)=-E(u,v),\qquad
 E_{ij}=0\quad\text{if }i+j\text{ is even}.
 \label{eq:parity}
\end{equation}
Thus \(E\) only couples even indices to odd indices.

Let \(S_{i,j}=\delta_{i,j+1}\).  Because \(T(-r/2)^{-1}=T(r/2)\),
restoring the boundary factors in \eqref{eq:B-kernel} gives
\begin{equation}
 A=MEM^{\mathsf T},\qquad
 M=(I+S)T(r/2).
 \label{eq:A-MEM}
\end{equation}
The matrix \(M\) is unit lower triangular.  Its relevant
first-subdiagonal entry is
\begin{equation}
 M_{2k,2k-1}
 =T(r/2)_{2k,2k-1}+T(r/2)_{2k-1,2k-1}
 =1+kr.
 \label{eq:M-subdiag}
\end{equation}

\section{Evaluation of the adjacent minor}

For any finite, even-cardinality row set \(I\), Pfaffian
minor summation gives
\begin{equation}
 \Pf(MEM^{\mathsf T})_I
 =\sum_{\substack{J\subset\Z_{\geq0}\\ |J|=|I|}}
 \det(M_{I,J})\,\Pf(E_J).
 \label{eq:minor-summation}
\end{equation}
Although the matrices have been written with infinite indices, the sum in
\eqref{eq:minor-summation} is finite: lower triangularity bounds every
column index by the largest row index.  Formula
\eqref{eq:minor-summation} may also be obtained directly by expanding the
alternating two-form associated with \(E\).

Set
\[
 I_0=\{0,\ldots,2k-1\},\qquad
 I_1=\{0,\ldots,2k-2,2k\}.
\]
For \(I_0\), unit lower triangularity leaves only \(J=I_0\), so
\begin{equation}
 C_0=\Pf E_{I_0}.
 \label{eq:C0-centred}
\end{equation}
For \(I_1\), the first \(2k-1\) rows force the columns
\(0,\ldots,2k-2\).  The last column is consequently either \(2k-1\)
or \(2k\).  Equations \eqref{eq:A-MEM} and
\eqref{eq:minor-summation} therefore give
\begin{equation}
 C_1=M_{2k,2k-1}\Pf E_{I_0}+\Pf E_{I_1}.
 \label{eq:C1-centred}
\end{equation}

The set \(I_1\) contains \(k+1\) even indices and \(k-1\) odd indices.
By \eqref{eq:parity}, every nonzero entry of \(E\) pairs indices of
opposite parity.  No perfect matching of \(I_1\) can therefore contribute,
and
\[
 \Pf E_{I_1}=0.
\]
Using \eqref{eq:M-subdiag} and \eqref{eq:C0-centred} in
\eqref{eq:C1-centred} proves \eqref{eq:minor-target}, without division by
\(C_0\).  Finally, Lemma~\ref{lem:reduction} gives
\begin{align*}
 [t^3]X^{\mathrm O}_{2n}(r,t)
 &=r^n\left((1-r-r^2)+r(1+(n-1)r)\right)C_0\\
 &=r^n\bigl(1+(n-2)r^2\bigr)C_0\\
 &=r\bigl(1+(n-2)r^2\bigr)X^{\mathrm O}_{2n-2}(r,1),
\end{align*}
which completes the proof of Theorem~\ref{thm:main}.

\begin{remark}
The Laurent extraction is initially performed over \(\Q(r)\), and the
binomial congruence over \(\Q[r]\).  The final equality is between
polynomials in \(\Z[r]\); it therefore includes \(r=0\) and does not depend
on dividing by a possibly vanishing partition function.
\end{remark}

\section{Finite-size interpretation and checks}

For positive \(r\), the coefficients of \eqref{eq:generator} are
nonnegative Boltzmann weights.  Theorem~\ref{thm:main} fixes the weighted
boundary one-point contribution at \(T=3\) in terms of the unrefined
partition function two sizes lower.  Coefficientwise,
\eqref{eq:coefficientwise} says that the weight-\(\rho\) sector receives
one shifted copy of the smaller model and \(n-2\) copies shifted by three
powers of \(r\).  This statement is exact at finite size and under the BFK
triangular/reflection boundary conditions.

At the first two sizes the identity reads
\[
 [t^3]X^{\mathrm O}_{4}(r,t)=r^2,\qquad
 [t^3]X^{\mathrm O}_{6}(r,t)=3r^3+3r^5,
\]
in agreement with direct definition-level enumeration and independent
evaluation of the source Pfaffian.  Exact integer-polynomial
implementations also agree coefficientwise through the reserved order-16
check.  These computations test conventions, signs and low-order cases;
the all-\(n\) conclusion follows from the parity argument above, not from
finite extrapolation.

No thermodynamic limit, boundary free energy, phase transition or
experimental observable is inferred here.  The result concerns the exact
finite partition polynomial and its formal weight.  The binomial-centred
form \eqref{eq:E-kernel} may nevertheless be useful for other nearby
boundary minors, since any centred index set with unequal even and odd
cardinalities has vanishing Pfaffian.

\section{Conclusion}

The formal-\(r\) third-boundary identity conjectured by BFK follows from a
short structural property of their Pfaffian kernel.  Binomial centring
exposes an opposite-parity support rule; Pfaffian minor summation then
reduces the shifted minor to a single first-subdiagonal entry of the
change-of-basis matrix.  This proves the conjecture for every \(n\geq2\)
while retaining the exact finite-size six-vertex boundary conditions.

\section*{Data and code availability}

The accompanying reproducibility package contains the exact symbolic
checker, two independent finite-size implementations, coefficient data and
the commands used to regenerate the reported checks.  No experimental or
personal data are involved.

\section*{Acknowledgements and generative-AI disclosure}

OpenAI Codex (model \texttt{gpt-5.6-sol}, reasoning effort
\texttt{ultra}) was used to assist with literature discovery, symbolic
derivation, exact-computation code, independent checking in a separate model
context, and generation and revision of the manuscript text.  Its outputs,
the cited sources and the underlying calculations require critical review by
the author before submission; the AI system is not an author and cannot
assume responsibility for the work.

\bibliographystyle{unsrtnat}
\bibliography{references}

\end{document}
